\documentclass[12pt]{article}
\usepackage[utf8]{inputenc}
\usepackage[T1]{fontenc}
\usepackage[margin=1in]{geometry}
\usepackage{setspace}
\usepackage{xcolor}
\usepackage[colorlinks=true, urlcolor=blue!60!black, linkcolor=blue!60!black]{hyperref}
\setlength{\parindent}{0em}
\setlength{\parskip}{0.6em}
\renewcommand{\baselinestretch}{1.0}
\begin{document}
\begin{center}
{\LARGE \textit{West Coast Hotel Co. v. Parrish}, 300 U.S. 379 (1937)\par}
\vspace{0.5em}
{\large Case Brief --- Two-Page Edition\par}
\end{center}

\section{Memory Jogger}
The ``switch in time'': the Court overruled \textit{Adkins} and upheld a minimum wage for women, ending \textit{Lochner}-era liberty of contract (\href{https://supreme.justia.com/cases/federal/us/300/379/}{300 U.S.\ 379}).

\section{Detailed Case Facts}
Washington's 1913 statute barred employing women at wages inadequate for their maintenance and let a state committee set binding minimums, here \$14.50 for a 48-hour week for hotel chambermaids. Elsie Parrish, a chambermaid at the West Coast Hotel Company's Cascadian Hotel, sued for \$216.19 in back wages; the hotel argued that the law violated freedom of contract under \textit{Adkins}.

\section{Procedural History}
The trial court struck the law under \textit{Adkins}; the Washington Supreme Court reversed. A year earlier, \textit{Morehead} had struck a similar New York law without reconsidering \textit{Adkins}.

\section{Judicial Votes}
Affirmed, 5--4. Hughes wrote for Brandeis, Stone, Cardozo, and Roberts; Sutherland dissented for Van Devanter, McReynolds, and Butler. Roberts voted in December 1936, before the court-packing plan was announced.

\section{Holding}
The Constitution guarantees no absolute freedom of contract. Regulation reasonably related to a proper purpose and not arbitrary or capricious satisfies due process, and a legislature may protect women first. \textit{Adkins} is overruled.

\section{Analysis of Opinions}
Hughes held that constitutional liberty is ``liberty in a social organization,'' subject to reasonable regulation, with ``every possible presumption'' favoring validity. He relied on the State's interest in women's health and bargaining power and observed that sub-living wages force taxpayers to provide ``a subsidy for unconscionable employers.'' Sutherland, dissenting, insisted that constitutional meaning does not change with economic conditions, called the law ``essentially arbitrary'' for ignoring the value of the work, and argued that women ``stand upon a legal and political equality with men.''

\section{Examples: Future Applications}
\textbf{Same-Side:} a gender-neutral \$15 minimum wage survives deferential review. \textbf{Opposite-Side:} a women-only wage law today fails intermediate scrutiny under \textit{United States v.\ Virginia}. \textbf{Fence-Sitter:} protectionist licensing of hair braiders.

\section{Critique}
Libertarians such as Barnett, Epstein, and Bernstein see an abdication that made economic liberty second class. Progressives praise the end of \textit{Lochner}, and Barry Cushman argues Roberts's vote reflected doctrine, not politics.

\section{Key Quotations}
\begin{quote}
``The Constitution does not speak of freedom of contract. It speaks of liberty and prohibits the deprivation of liberty without due process of law. In prohibiting that deprivation, the Constitution does not recognize an absolute and uncontrollable liberty. Liberty in each of its phases has its history and connotation. But the liberty safeguarded is liberty in a social organization which requires the protection of law against the evils which menace the health, safety, morals, and welfare of the people. Liberty under the Constitution is thus necessarily subject to the restraints of due process, and regulation which is reasonable in relation to its subject and is adopted in the interests of the community is due process.''
\end{quote}
\noindent --- Chief Justice Hughes, Majority Opinion, 300 U.S.\ at 391.

\end{document}
